\documentclass{article}

\textwidth=6.5in
\textheight=8.5in
\voffset=-54pt
\oddsidemargin=0in
\evensidemargin=0in
\setlength{\parskip}{8pt}
\setlength{\parindent}{0pt}

\usepackage{amsmath, amssymb}
\usepackage{ifthen}

\usepackage[inline]{enumitem}
\setlist{topsep=0pt, parsep=8pt}

\usepackage[rgb,table]{xcolor}\convertcolorsUfalse
\usepackage{graphicx}

% change hyperref options
\PassOptionsToPackage{hyphens}{url}
\usepackage[bookmarks,colorlinks,linkcolor=black,citecolor=blue,pdfstartview=FitH,urlcolor=blue]{hyperref}
\hypersetup{pdfpagemode=UseNone}
\usepackage{bookmark}

\usepackage{graphicx}
\usepackage{multicol}
\usepackage{framed}
\usepackage{array}

\usepackage{icomma}

\usepackage{soul}

\usepackage{tikz}
\usetikzlibrary{
  matrix,%
  % enables the snake paths
  decorations.pathmorphing,%
  % enables bigger arrows
  decorations.markings,%
  decorations.pathreplacing,%
  decorations.text,%
  calc,%
  patterns,%
  shapes.geometric,%
  arrows,
  decorations.shapes,
  positioning,
  plotmarks,
  intersections
}
\tikzset{>=latex} % sets default arrow type in tikz
\tikzset{font=\normalsize}
\tikzset{mark size=1.5pt, mark options=thin}
\tikzset{pin distance=4pt,
  every pin edge/.style={<-, thin, shorten <= -2pt}}

\interfootnotelinepenalty=10000
\renewcommand{\thefootnote}{(\arabic{footnote})}

\usepackage{multirow, bigdelim}

\usepackage[notref,notcite]{showkeys}

% change to \usepackage[sols]{handout} to see solutions
\usepackage[sols,notes,workshop,grading]{handout}

\newcommand\iscurrentenumi[1]{%
  \ifthenelse{\equal{#1}{\theenumi}}{}{#1}}
\setenumerate[1]{ref=\#\arabic*}
\setenumerate[2]{ref=\protect\iscurrentenumi{\theenumi}(\alph*)}
\newcommand\enumiiprefix[2]{%
  \ifthenelse{{\equal{#1}{\theenumi}}\AND{\equal{#2}{(\alph{enumii})}}}{}{}%
  \ifthenelse{{\equal{#1}{\theenumi}}\AND{\NOT{\equal{#2}{(\alph{enumii})}}}}{#2}{}%
  \ifthenelse{\equal{#1}{\theenumi}}{}{#1#2}%
}
\setenumerate[3]{ref=\protect\enumiiprefix{\theenumi}{(\alph{enumii})}\roman*}

\definecolor{purple}{rgb}{0.5,0,1.0}

\usepackage{fancyhdr}
\lhead{\textcolor{white}{This content is protected and may not be shared, uploaded, or distributed.}}
\rhead{}
\rfoot{\vspace*{\fill}\footnotesize\textcopyright{} \emph{Harvard Math Ma}}
\renewcommand{\headrulewidth}{0pt}
\pagestyle{fancy}
\begin{document}

\begin{center}
  \Large\textsc{Learning Objectives}
\end{center}

\begin{enumerate}
\item \textbf{Introduction to Functions} %% 1
  \begin{itemize}

  \item Explain what a \emph{function} is, and distinguish between functions and non-functions.

  \item Identify the \emph{independent variable}, \emph{dependent variable}, \emph{domain}, and \emph{range} of a function.

  \item Explain how slope graphically represents a rate of change.

  \item Reason about whether a function is \emph{increasing} or \emph{decreasing}, as well as whether it's \emph{concave up}, \emph{concave down}, or \emph{linear}.

  \end{itemize}

\item \textbf{A Library of Basic Functions and Properties} %% 2
  \begin{itemize}
  \item Sketch and recognize graphs of $y = mx + b$, $y = x^2$, $y = x^3$, $y = \frac{1}{x}$, $y = \frac{1}{x^2}$, $y = |x|$, and $y = \sqrt{x}$.

  \item Understand the following characteristics of a function graphically: positive / negative, increasing / decreasing, intercepts, asymptotes, continuity, concavity, even / odd.

  \item Algebraically determine whether a function is even, odd, or neither.

  \item Explain what it means to say a function is increasing / decreasing \emph{without bound}, as opposed to having an \emph{upper bound} or \emph{lower bound}.

  \item Recognize when the graph of a function given algebraically has a hole and construct algebraic examples of functions whose graphs have holes.

  \end{itemize}

\item \textbf{More on Functions, Average Rates of Change} %% 3
  \begin{itemize}

  \item Write and interpret expressions in function notation.

  \item Compute an average rate of change, interpret it in the context of a situation, and represent it graphically.

  \item Explain the difference between average speed, average velocity, instantaneous speed, and instantaneous velocity.

  \item Explain how concavity can be characterized in terms of secant line segments.

  \end{itemize}

\item \textbf{Transforming Functions} %% 4
  \begin{itemize}
  \item Sketch $y = a f(bx) + c$, $y = a f(x + b) + c$, and $y = |f(x)|$
    from a graph of $y = f(x)$.

  \end{itemize}

\item \textbf{Linear and Piecewise Linear Functions} %% 5
  \begin{itemize} 
  \item Write the equation of a line given 2 points or 1 point and the
    slope, and decide whether slope-intercept form ($y = mx + b$) or
    point-slope form is more useful. 

  \item Use linear and piecewise linear equations to model various
    phenomena.

  \item Interpret the slope in linear and piecewise linear functions, 
    and graph the ``slope function'' $f'(x)$ of a piecewise linear
    function $f(x)$.  

  \end{itemize}

\item \textbf{Function Composition, Modeling with Functions} %% 6
  \begin{itemize}
  \item Explain how composition of functions models situations where there is a chain of dependency among variables.

  \item Compute compositions of functions.

  \item Model situations systematically: identify your goal, define variables and express your goal in terms of those variables, and then look for relationships among your variables.

  \item Communicate your modeling process clearly. 
  \end{itemize}

  \addtocounter{enumi}{1}
  
%\item \textbf{Problem-solving day} %% 7
\item \textbf{Approximating with Secant Lines} %% 8
  \begin{itemize}
  \item Decide where the graph of a function is or isn't locally linear.

  \item Use secant lines to approximate function values.

  \item Determine whether a secant line approximation is too high or too low, or if there isn't enough information to be sure.

  \end{itemize}

\item \textbf{Definition and Interpretation of the Derivative} %% 9
  \begin{itemize}
  \item State the definition of the derivative, explain where it comes from, and use it to calculate derivatives.

  \item Interpret the derivative in context and represent it graphically.

  \item Find the tangent line to the graph of $f(x)$ at $x = a$. 

  \item Use tangent lines to approximate function values, and determine whether such an approximation is too high or too low, or if there isn't enough information to be sure.
  \end{itemize}

\item \textbf{The Derivative Function} %% 10
  \begin{itemize}
  \item Use the limit definition to compute the derivative function
    $f'$.

  \item Determine whether a function is \emph{differentiable}, both
    graphically and algebraically, and give examples of functions that
    are differentiable and non-differentiable.

  \item Sketch a rough graph of $f'$ when you have the graph of $f$. 

  \end{itemize}

\item \textbf{Functions and their Derivatives} %% 11
  \begin{itemize}

  \item Explain how concavity can be characterized in terms of tangent lines.

  \item Explain what $f'$ and $f''$ tell us about the graph of $f$. (For example, if $f'$ is positive, what does that say about $f$? If $f'$ is increasing, what does that say about $f$?)

  \item Explain how displacement, velocity, and acceleration are related in terms of derivatives.

  \end{itemize}

\item \textbf{Using Derivatives to Graph Functions} %% 12
  \begin{itemize}

  \item Use the first derivative to determine where a function is
    increasing or decreasing and to find \emph{turning points}.

  \item Use the second derivative to determine where a function is
    concave up or concave down and to find \emph{inflection points}. 

  \item Put the above information together to sketch a reasonable graph
    of a function.

  \item Develop the habit of checking for consistency when putting several pieces of information together. 
  \end{itemize}

\item \textbf{Limits} %% 13
  \begin{itemize}
  \item Calculate two-sided and one-sided limits of functions given algebraically or graphically, and understand why such a limit might not exist.

  \item Recognize when we need to factor and cancel to calculate a limit (for example, with $\displaystyle \lim_{x \to 2} \frac{x^2 - 4}{x - 2}$).

  \item Explain precisely what we mean when we write $\displaystyle \lim_{x \to a} f(x) = \infty$ or $\displaystyle \lim_{x \to a} f(x) = -\infty$. In particular, you should understand that, in both of these cases, $\displaystyle \lim_{x \to a} f(x)$ does not exist, but that saying $\displaystyle \lim_{x \to a} f(x) = \infty$ or $\displaystyle \lim_{x \to a} f(x) = -\infty$ gives more information than just saying ``$\displaystyle \lim_{x \to a} f(x)$ does not exist'' does.

  \item Use algebraic methods to determine whether a limit is $\pm \infty$.

  \end{itemize}

\item \textbf{More on Limits} %% 14
  \begin{itemize}
  \item Evaluate $\displaystyle \lim_{x \to \infty} f(x)$ and $\displaystyle \lim_{x \to -\infty} f(x)$ and interpret such limits graphically.

  \item Find the horizontal asymptotes of a function $f(x)$ if given an algebraic formula for $f(x)$. 
  \end{itemize}
  
%\item \textbf{Exam 1 Practice} %% 15

    \addtocounter{enumi}{1}

\item \textbf{Continuity} %% 16
  \begin{itemize}
  \item Explain the idea of continuity, both intuitively and precisely using the limit definition.

  \item Use limits to classify discontinuities as \emph{removable discontinuities}, \emph{jump discontinuities}, or \emph{vertical asymptotes}, and describe what each type of discontinuity looks like graphically.

  \item Explain what relationships there are between whether a function is continuous at $a$ and whether it's differentiable at $a$.

  \item Take discontinuities into account when constructing a sign chart for a function. 
  \end{itemize}

\item \textbf{3 Differentiation Rules} %% 17
  \begin{itemize}
  \item Explain why the Power Rule is true for positive integer exponents.

  \item Explain why the Constant Multiple Rule and Sum Rule are true.

  \item Use the Power Rule, Constant Multiple Rule, and Sum Rule, and recognize when they do and do not apply.

  \end{itemize}

\item \textbf{More Differentiation Rules, Tangent Line Approximation} %% 18
  \begin{itemize}
  \item Explain why the Product and Quotient Rule are true, and use them (together with the rules we've learned before) to differentiate functions.

  \item Identify functions that can and can't be differentiated using just the rules we've learned so far.

  \item Use secant lines and tangent lines to approximate values like $\sqrt{22}$, and use concavity to determine whether your approximations are upper or lower bounds. 

  \end{itemize}
  \item \textbf{The Chain Rule} %% 18
  \begin{itemize}
  \item You should be able to use the Chain Rule to find the derivative of
  the composition of two functions.

\item You should understand why the Chain Rule makes sense in a real-world
  problem (such as {{\#1}} on {the worksheet ``The Chain Rule''}).

\item You should be comfortable looking at a function and determining its
  structure: is it a product?  A composition?  (If so, what is the outer
  function?)

  \end{itemize}

\item \textbf{Maxima and Minima of Functions} %% 19
  \begin{itemize}

  \item Explain the difference between a \emph{local minimum / maximum} and a \emph{global (or absolute) minimum / maximum}.

  \item Find the local extrema of a function by first finding \emph{critical points} and then using the First Derivative Test or Second Derivative Test.

  \item Explain why the First and Second Derivative Tests work.

  \end{itemize}

\item \textbf{Finding Global Extrema} %% 20
  \begin{itemize}
  \item Explain what the Extreme Value Theorem says.

  \item Determine whether a given function has global extrema on a given domain and, if so, find the global extrema.

  \end{itemize}

\item \textbf{Optimization} %% 21
  \begin{itemize}
  \item Set up and solve optimization problems using calculus.

  \item Justify that your solution is optimal using appropriate techniques from calculus.

  \end{itemize}

\item \textbf{Introduction to Exponential Functions} %% 22
  \begin{itemize}
  \item Model quantities that have a constant percent change (per unit of input). 

  \item Graph $b^x$ for different values of $b$ (and be able to combine this with what we know about altering functions to graph functions involving $b^x$).

  \item Apply exponent rules. 
  \end{itemize}

\item \textbf{Applications of Exponentials} %% 23
  \begin{itemize}
  \item Model constant percent change over different intervals (like ``decreases by 30\% every 12 hours'' or ``triples every 4 days''), including compound interest.

  \item Explain what \emph{half-life} means, and model phenomena that have a half-life. 

  \item Reason about how an exponential function grows or decays over different time intervals. (For example, if a quantity grows by 10\% every year, will it grow by exactly 20\%, more than 20\%, or less than 20\% in 2 years?)

  \end{itemize}

\item \textbf{Derivatives of Exponential Functions} %% 24
  \begin{itemize}
  \item Explain what we mean when we say ``the derivative of an exponential function is proportional to the function itself''.

  \item Explain how the number $e$ is defined in terms of derivatives.

  \item Differentiate $e^{kx}$ where $k$ is any constant.

  \end{itemize}

\item \textbf{More on Exponentials} %% 25
  \begin{itemize}
  \item Fit an exponential function to two data points, a graph, or a verbal description.

  \item Apply the fact that the derivative of an exponential function is proportional to the function.

  \end{itemize}

%\item \textbf{Problem-Solving Day} %% 26
%\item \textbf{Exam 2 Practice} %% 27

  \addtocounter{enumi}{1}

\item \textbf{Inverse Functions} %% 28
  \begin{itemize}

  \item Determine whether a function $f$ is invertible based on a graph,
    description, or formula. If so, determine the domain and range of $f^{-1}$. 

  \item If $f$ is an invertible function given verbally (for example, if a farmer is growing a giant pumpkin and $f(t)$ gives the pumpkin's weight at time $t$), describe what $f^{-1}$ represents. 

  \item Find a formula for $f^{-1}$ when given a formula for an invertible function $f$.

  \item Sketch the graph of $f^{-1}$ from a graph of an invertible function $f$.

  \end{itemize}

\item \textbf{Introduction to Logarithms} %% 29
  \begin{itemize}
  \item State and apply the definition of a logarithm, and explain how logarithms are related to exponentials.

  \item Explain how to simplify $\log_b \left( b^{\bigstar} \right)$ and $b^{\log_b \heartsuit}$.

  \item Sketch $y = \log_b x$ and explain how the graph and its important features (domain, range, asymptotes, intercepts, etc.)\ relate to the graph of $y = b^x$.

  \item Differentiate $b^x$ for any constant $b > 0$.

  \end{itemize}
\addtocounter{enumi}{1}
\item \textbf{Logarithm Rules} %% 30
  \begin{itemize}

  \item Apply basic properties of logarithms. (Can we simplify the log of a product? How about the log of a sum, quotient, or power?) 

  \item Change bases of logarithms (for example, rewrite $\log_2 3$ in terms of $\ln$ or $\log_7$). 

  \end{itemize}

\item \textbf{Derivatives of Logarithmic and Exponential Functions} %% 31
  \begin{itemize}
  \item Differentiate logarithmic functions, and apply this together with the differentation rules we've learned before. 

  \item Use algebra to rewrite functions involving exponentials or logarithms so that they can be differentiated.

  \end{itemize}
\item \textbf{The Chain Rule, Continued} %% 32
  \begin{itemize}
 \item You should be able to apply the Chain Rule like a pro!

\item You should be able to apply the Chain Rule in more general contexts,
  such as in a word problem or when given a composition of functions
  described by graphs or numerical data.

\item You should understand how logarithms and the Chain Rule can be used
  to show that the Power Rule $\tfrac{d}{dx} \left( x^n \right) =
  nx^{n-1}$ holds for any real number $n$. 


  \end{itemize}
\item \textbf{Solving Equations with Logarithms and Exponents} %% 32
  \begin{itemize}
  \item Solve equations involving exponentials and logarithms.

  \end{itemize}

\item \textbf{Linear Functions, Exponentials, and Logarithms in the Real World} %% 33
  \begin{itemize}
  \item Assess how appropriate it is to use a linear or exponential model to model data.

  \item Interpret plots where an axis has a logarithmic scale. 

  \end{itemize}


\item \textbf{Cubics and Higher-Degree Polynomials} %% 35
  \begin{itemize}
  \item Identify the degree of a polynomial, and explain what this tells us about the possible number of roots and turning points of the polynomial. 

  \item Describe the ``end behavior'' of a polynomial, and explain what this tells us about the possible number of roots and turning points of the polynomial.

  \item Write a formula for a polynomial if given enough information about its characteristics (for example, roots, local extrema, etc.).

  \end{itemize}

\item \textbf{Comparing Exponentials, Logarithms, and Power Functions} %% 36
  \begin{itemize}
  \item Describe the relative growth rates of logarithmic functions, power functions (functions of the form $f(x) = x^p$ where $p$ is a constant), and exponential functions, and use these to calculate limits.    

  \end{itemize}

\end{enumerate}
\end{document}
